\section{Assertions}
In this project, \textbf{assertions} are used as a tool to catch programming errors and invalid states during development. They help ensure that assumptions made in the code are correct, and they provide valuable debugging information when things go wrong. However, assertions should never be relied on to handle runtime errors or expected conditions in production code.

\subsection{Guidelines}
\begin{itemize}
    \item \textbf{Use Assertions for Debugging}: Assertions are intended to catch conditions that should never occur during normal program execution. Use them to assert assumptions about the state of the program, such as parameter validation or invariants within algorithms.
\end{itemize}
\begin{lstlisting}[language=C]
#include <assert.h>

void process_data(int* data)
{
  assert(data != NULL);  /* Assert that data is not NULL */
  /* Continue processing */
}
\end{lstlisting}

\subsection{ESP32-Specific Assertions}
In ESP32 projects, you have two specific macros for error checking: \textbf{\texttt{ESP\_ERROR\_CHECK}} and \textbf{\texttt{ESP\_ERROR\_CHECK\_WITHOUT\_ABORT}}. These are commonly used to catch errors during development, especially when working with ESP-IDF functions.

\begin{itemize}
    \item \textbf{\texttt{ESP\_ERROR\_CHECK}}: This macro checks the return value of an ESP-IDF function and aborts the program if the return code is not \texttt{ESP\_OK}.
    \item \textbf{\texttt{ESP\_ERROR\_CHECK\_WITHOUT\_ABORT}}: This macro works similarly to \texttt{ESP\_ERROR\_CHECK}, but it will not abort the program when an error occurs.
\end{itemize}
